% !PRINTED 293
\weberpage{293}
Wir fügen nun zu den beiden Functionen $\varphi, \psi$ eine dritte Function $f$ hinzu, die, gleich Null gesetzt, eine \gesperrt{geschlossene} Curve darstellt, und bestimmen die Charakteristik des Functionensystems ganz in der früheren Weise, indem wir längs der Curve $f$ fortschreiten
\weq{1}{k = \tfrac{1}{2}\big[E(f;\varphi,\psi) - A(f;\varphi,\psi)\big] = \tfrac{1}{2}\big[E(f;\psi,\varphi) - A(f;\psi,\varphi)\big].}
So ist z.\,B. in der Fig. 11 die Charakteristik $k = 2$.

\weberfigureimage{weber-b1-fig11-12.png}{0.96}

Um unsere Sätze anwenden zu können, müssen wir die Curven $\varphi, \psi$ irgendwie durch willkürlich hergestellte Verbindungen, die aber alle ausserhalb $f$ verlaufen sollen, zu geschlossenen machen. Die Charakteristik dieses Systems ist dann gleichfalls durch (1) bestimmt. Als Beispiel ergänzen wir die Fig. 11 durch Fig. 12, in der diese abschliessenden Verbindungen gezeichnet sind. Die so hinzugefügten Schnittpunkte können theils Austrittspunkte $A(\varphi,\psi)$, theils Eintrittspunkte $E(\varphi,\psi)$ sein. Wir wollen ihre Anzahlen mit
\[
A'(\varphi,\psi), \quad E'(\varphi,\psi)
\]
bezeichnen, während $A(\varphi,\psi)$ die Zahl der ursprünglich vorhandenen Schnittpunkte bedeutet. Immer muss jetzt
\weq{2}{A(\varphi,\psi) + A'(\varphi,\psi) = E'(\varphi,\psi)}
sein. Wir theilen die Punkte $A(\varphi,\psi)$ in zwei Gruppen $A_\alpha(\varphi,\psi)$, $A_\varepsilon(\varphi,\psi)$, von denen die ersten ausserhalb, die anderen innerhalb von $f$ liegen sollen. Dann haben wir folgende Charaktere der Schnittpunkte
%CONT>>
